\section{Strategy: Adjoint Interpolation and Center Symmetry}
\label{sec:adjoint}

\subsection{Motivation}
A central difficulty in proving the mass gap for Pure Yang-Mills theory is the lack of a perturbative starting point that exhibits a gap. The weak coupling limit is gapless in perturbation theory, and the strong coupling limit, while gapped, is far from the continuum.

We discuss a strategy based on interpolating between Pure Yang-Mills theory and $\mathcal{N}=1$ Supersymmetric Yang-Mills (SYM) theory, where a mass gap is known to exist.

\subsection{The Interpolating Action}
We consider a lattice action $S(M)$ that couples the gauge field to a massive adjoint Majorana fermion (gluino) with mass $M$:
\begin{equation}
S(M) = S_{YM} + S_{Majorana}(M)
\end{equation}
The parameter $M$ serves as an interpolating parameter:
\begin{itemize}
    \item As $M \to 0$, we recover $\mathcal{N}=1$ SYM.
    \item As $M \to \infty$, the fermions decouple, and the effective theory approaches Pure Yang-Mills.
\end{itemize}

\subsection{Properties at the Endpoints}
\subsubsection{The SUSY Limit ($M=0$)}
In the $M=0$ limit, the Witten Index $\text{Tr}(-1)^F$ can be computed. For periodic boundary conditions, $I_W \neq 0$ for $SU(N)$, which implies that supersymmetry is unbroken. Standard arguments in supersymmetric field theories suggest that unbroken supersymmetry in 4D implies a mass gap, provided there are no massless Goldstone modes. The breaking of the $U(1)_R$ symmetry to a discrete subgroup suggests no Goldstone bosons appear.

\subsubsection{The Pure YM Limit ($M \to \infty$)}
This is the target theory.

\subsection{The Continuity Hypothesis}
The critical question is whether the mass gap remains open for all $M \in [0, \infty)$.
Both endpoints share the same realization of Center Symmetry ($Z_N$).
\begin{enumerate}
    \item At $M=0$, the theory is expected to be confining ($\langle P \rangle = 0$).
    \item At $M \to \infty$, the theory is confining.
\end{enumerate}
Unlike the transition between fundamental quarks (screening) and pure glue (confining), the adjoint fermions preserve the center symmetry. Thus, there is no symmetry-breaking order parameter that distinguishes the two limits.

\begin{conjecture}[Adiabatic Continuity]
There exists a path in the $(g^2, M)$ plane connecting the $M=0$ regime to the $M \to \infty$ regime along which no phase transition (bulk or boundary) occurs.
\end{conjecture}

If this conjecture holds, then the mass gap established at $M=0$ would persist to the pure Yang-Mills theory. However, we emphasize that the absence of a first-order phase transition (e.g., a liquid-gas type transition) is not guaranteed by symmetry alone and constitutes a major open problem.

\subsection{Relationship to Previous Approaches}
Prior attempts to prove this continuity have relied on assumptions about the analyticity of the free energy (e.g., using Pirogov-Sinai theory). Proving such analyticity uniformly in the volume is equivalent to solving the gap problem itself. We therefore state this Continuity as a necessary condition for this strategy, rather than a proven theorem.
